% Corrected French for printed p.162.
% Correction: sections (plural).
\SourcePage{162}

\LayerTarget{tohoku:prop:3.4.2:p2}
\textit{suite exacte de faisceaux de groupes sur $X$, avec $F$ abélien.
Soient $E$ un espace fibré sur $X$ à faisceau structural $H$ [11, Chap. IV],
$F^E$ le faisceau de groupes associé à $E$ et aux opérations de $H$ sur $F$
définies à l'aide des automorphismes intérieurs de $G$ (qui opèrent sur le
faisceau invariant $F$). On peut alors définir un \og élément cobord\fg{}
$\partial E\in H^2(X,F^E)$, de façon \og fonctorielle\fg{}, de telle façon que
la condition nécessaire et suffisante pour que $\partial E=0$, c'est que la
classe $c(E)$ de $E$ dans $H^1(X,H)$ [11, Chap. V] appartienne à l'image de
$H^1(X,G)$.}

Cet énoncé se simplifie quand $F$ est dans le centre de $G$, car alors
$F^E=F$ ne dépend plus du fibré $E$, et on obtient :

\LayerTarget{tohoku:cor:3.4.2}
\textsc{Corollaire.} --- \textit{Soit
$e\longrightarrow F\longrightarrow G\longrightarrow H\longrightarrow e$
une suite exacte de faisceaux de groupes sur l'espace $X$, avec $F$ dans le
centre de $G$. Alors on peut trouver une application (\og fonctorielle\fg{})
\[
\partial:H^1(X,H)\longrightarrow H^2(X,F)
\]
telle que $\partial^{-1}(0)=\operatorname{Im}(H^1(X,G))$.}

\LayerTarget{tohoku:proof:3.4.2:p1}
\textsc{Démonstration de prop. 3.4.2.} Plongeons tout faisceau $M$ sur $X$
dans le faisceau $\overline M$ dont l'ensemble des sections sur l'ouvert $U$
est l'ensemble

\LayerTarget{tohoku:construction:flabby-envelope}
\[
\prod_{x\in U}M(x);
\]
$\overline M$ est donc un faisceau \textit{flasque} et $M$ s'identifie à un
sous-faisceau de $\overline M$. Si $M$ est un faisceau de groupes (resp. de
groupes abéliens) il en est de même de $\overline M$. D'autre part, on vérifie
facilement que si $L$ est un faisceau flasque de groupes (non nécessairement
abéliens), alors $H^1(X,L)$ est réduit à l'élément neutre (en d'autres termes,
tout faisceau \textit{principal} sous $L$ [11, définition 3.4.2] admet une
section ; on construit aisément une telle section par \og zornification\fg{}
sur l'ensemble des sections construites sur des ouverts de $X$).

\LayerTarget{tohoku:comparison:cech-derived-h1}
Il résulte de ceci, que pour des faisceaux $F$ de groupes abéliens, le groupe
$H^1(X,F)$ tel qu'il est défini dans [11] (par la méthode de Čech) est le même
que celui défini axiomatiquement dans ce travail (savoir le premier satellite
droit $S^1\Gamma$ du foncteur $\Gamma$ sur $C^X$). Revenons alors aux
conditions de prop. 3.4.2, on a un homomorphisme de suites exactes :

\LayerTarget{tohoku:diagram:3.4.2:flabby-embedding}
\[
\begin{array}{ccccccccc}
e&\longrightarrow&F&\longrightarrow&G&\longrightarrow&H&\longrightarrow&e\\
&&\big\downarrow&&\big\downarrow&&\big\downarrow&&\\[0.45em]
e&\longrightarrow&\overline F&\longrightarrow&\overline G&\longrightarrow
&\overline H&\longrightarrow&e
\end{array}
\qquad \overline F\cap G=F
\]
$G$ et $\overline F$ sont des sous-faisceaux de groupes de $\overline G$,
$\overline F$ distingué. Soit $\overline F\mathbin{\cdot}G$ le faisceau de
groupes de $\overline G$ engendré par $G$ et $\overline F$, posons :

\LayerTarget{tohoku:construction:3.4.2:P-Fprime}
\LayerTarget{tohoku:eq:3.4.2:quotients}
\[
P=\overline F G,\qquad F'=\overline F/F.
\]
On définit de façon évidente une suite exacte d'homomorphismes de faisceaux de
groupes

\LayerTarget{tohoku:eq:3.4.2:P-extension}
\[
e\longrightarrow\overline F\longrightarrow P\longrightarrow H
\longrightarrow e.
\]

\LayerTarget{tohoku:claim:3.4.2:H1-bijection}
Montrons que l'application correspondante
$H^1(X,P)\longrightarrow H^1(X,H)$ est bijective. C'est un monomorphisme en
vertu de la suite exacte de cohomologie de [11, prop. 5.6.2], compte tenu du
fait que $H^1(X,L)$ est réduit à l'élément neutre

\TohokuContinuation{163}
